\section{Introduction}

In four years, three in four Albanian fifteen-year-olds crossed the line into low
mathematics proficiency. The OECD Programme for International Student Assessment
(PISA) measures whether students can apply reading, mathematics, and science to real
problems, and its \emph{Level~2} boundary marks basic proficiency; students below it
are the international ``low performers'' that education systems target. In PISA~2018,
42\% of Albanian students fell below the mathematics Level~2 line. In PISA~2022, 75\%
did. No other participating system regressed so far in a single cycle. A drop of this
magnitude invites an obvious reading: the 2022 cohort was simply weaker, and the
country needs broad remediation. Whether that reading is right is not a semantic
question. It determines whether policy should treat 2022 as more of the same, only
worse, or as evidence that \emph{who} is at risk, and \emph{why}, has changed.

This paper treats the question as one of educational data mining, and it insists on
the survey methodology that PISA requires. PISA is a stratified, clustered sample
reported through final student weights, and each achievement score arrives as ten
\emph{plausible values}, multiple imputations of a latent proficiency. Analyses that
ignore the weights do not estimate a population quantity, and analyses that collapse
the plausible values to a single number understate their uncertainty. A large body of
secondary PISA work nonetheless does both. Every estimate below is survey-weighted;
plausible values are combined by Rubin's rules; and where the data permit (the
2015, 2018, and 2022 cycles carry 80 balanced-repeated-replication weights) standard
errors are design-based.

From the student background questionnaire, socioeconomic status, home possessions,
parental education and occupation, grade repetition, immigration status, school
belonging, teacher support, math anxiety, and access to information technology, we ask
three questions. Can a classifier identify Albanian students at risk of low
mathematics proficiency, and how well? Which factors drive that risk, and do they
differ from the drivers in comparable countries? And does the 2022 collapse represent
a change in the risk \emph{structure} that can explain the 2018-to-2022 reversal? The
economic and structural conditions of Albania in this exact window, a destructive 2019
earthquake, an extended pandemic school closure, a thin remote-learning infrastructure,
and sustained emigration of working-age adults and teachers, give these questions a
concrete backdrop that we return to in Section~\ref{sec:context}.

\paragraph{Thesis.}
Albania's 2022 mathematics collapse is not a uniformly weaker cohort but a
\emph{structural break in the risk-generating process}, driven by a disruption to the
supply of schooling rather than to household resources; the students at risk are
identifiable from background data only up to a genuine ceiling of about 0.78 AUC that
three independent methods confirm, their risk is broad-based rather than concentrated
in the poor, and its drivers are the same ones that operate across the region.

\paragraph{Roadmap.}
Section~\ref{sec:data} describes the five-cycle data and the design-based estimation.
Section~\ref{sec:methods} sets out the leakage-safe, survey-weighted modelling and the
covariate-shift diagnostic. Section~\ref{sec:results} reports the four empirical
findings: the structural shift (AUC~0.98), the feature-versus-data ceiling, the shared
and broad-based driver structure, and the calibrated, fairness-audited screener.
Section~\ref{sec:context} then reads those findings against Albania's economic and
structural conditions, arguing that a supply-side shock, not a household-income shock,
best fits the evidence. Section~\ref{sec:discussion} discusses implications and limits,
and Section~\ref{sec:conclusion} returns to the collapse with which we began.

\paragraph{Contributions.}
Our central contribution is methodological and, we believe, transferable beyond this
dataset: we show that the accuracy ceiling recurring across PISA prediction studies is
not one phenomenon but two, a movable \emph{feature} ceiling and a fixed \emph{data}
ceiling, and that the two can be told apart with convergent evidence rather than with
another round of tuning. Applied to Albania, this lets us (i) show the 2022 collapse is
a structural covariate shift rather than a level change; (ii) locate the background-data
accuracy ceiling, demonstrate it is a feature ceiling by moving it with school
socioeconomic composition, and then certify the new plateau as a genuine data ceiling
through three independent routes; (iii) characterise the crisis as broad-based and
shared with regional peers; and (iv) connect these statistical findings to Albania's
economic and structural context, while delivering a calibrated, interpretable, and
honestly-bounded risk screener that we argue should not be read causally.
